\section{Exact Sequences}

\begin{definition}
    Let $\{M_i\}_{i\in \mathbb Z}$ be a family of $R$-modules, and let
    \[
        \varphi_i:M_i\to M_{i+1}
    \]
    be $R$-module homomorphisms for each $i\in\mathbb Z$. Then
    \[
        \cdots \longrightarrow M_{i-1}\xrightarrow{\varphi_{i-1}} M_i \xrightarrow{\varphi_i} M_{i+1}\longrightarrow \cdots
    \]
    is called a \emph{sequence} of $R$-modules.
\end{definition}

\begin{definition}
    The sequence
    \[
        \cdots \longrightarrow M_{i-1}\xrightarrow{\varphi_{i-1}} M_i \xrightarrow{\varphi_i} M_{i+1}\longrightarrow \cdots
    \]
    is said to be \emph{exact at $M_i$} if
    \[
        \operatorname{Im}(\varphi_{i-1})=\ker(\varphi_i).
    \]
    It is called an \emph{exact sequence} if it is exact at every term.
\end{definition}

A short exact sequence is a sequence of the form
\[
    0\longrightarrow M' \xrightarrow{\iota} M \xrightarrow{\pi} M'' \longrightarrow 0.
\]

It is exact if and only if
\[
    \ker(\iota)=0,
    \qquad
    \operatorname{Im}(\pi)=M'',
    \qquad
    \operatorname{Im}(\iota)=\ker(\pi).
\]
Equivalently,$\iota$ is injective,$\pi$ is surjective, and
\[
    \operatorname{Im}(\iota)=\ker(\pi).
\]

Hence, for every short exact sequence
\[
    0\longrightarrow M' \xrightarrow{\iota} M \xrightarrow{\pi} M'' \longrightarrow 0,
\]
we have
\[
    M/M' \cong M''.
\]

\begin{definition}
    Let $R,S$ be rings, and let
    \[
        F: R\text{-}\mathrm{Mod} \to S\text{-}\mathrm{Mod}
    \]
    be a covariant functor.

    We say that $F$ is \emph{left exact} if for any exact sequence
    \[
        0 \to N \xrightarrow{\varphi} M \xrightarrow{\psi} P,
    \]
    the sequence
    \[
        0 \to F(N) \xrightarrow{F(\varphi)} F(M) \xrightarrow{F(\psi)} F(P)
    \]
    is exact.

    We say that $F$ is \emph{right exact} if for any exact sequence
    \[
        N \xrightarrow{\varphi} M \xrightarrow{\psi} P \to 0,
    \]
    the sequence
    \[
        F(N) \xrightarrow{F(\varphi)} F(M) \xrightarrow{F(\psi)} F(P) \to 0
    \]
    is exact.
\end{definition}

\begin{definition}
    Let
    \[
        G: R\text{-}\mathrm{Mod} \to S\text{-}\mathrm{Mod}
    \]
    be a contravariant functor. We say that $G$ is \emph{left exact} (respectively, \emph{right exact}) if it is left exact (respectively, right exact) when regarded as a covariant functor
    \[
        G : (R\text{-}\mathrm{Mod})^{\mathrm{op}} \to S\text{-}\mathrm{Mod}.
    \]

    A functor is called \emph{exact} if it is both left exact and right exact.
\end{definition}

Let $M$ be an $R$-module. Then:
\[
    \operatorname{Hom}_R(M,-): R\text{-}\mathrm{Mod} \to \mathrm{Ab}
    \quad\text{is left exact,}
\]
\[
    \operatorname{Hom}_R(-,M): R\text{-}\mathrm{Mod} \to \mathrm{Ab}
    \quad\text{is left exact,}
\]
and
\[
    M \otimes_R - : R\text{-}\mathrm{Mod} \to \mathrm{Ab}
    \quad\text{is right exact.}
\]

\begin{definition}
    A functor $F: \mathcal C \to \mathcal D$ is called \emph{exact} if it is both left exact and right exact.
\end{definition}

\begin{definition}
    Let $M$ be an $R$-module.

    \begin{enumerate}
        \item $M$ is called \emph{projective} if the functor
              \[
                  \operatorname{Hom}_R(M,-)
              \]
              is exact.

        \item $M$ is called \emph{injective} if the functor
              \[
                  \operatorname{Hom}_R(-,M)
              \]
              is exact.

        \item $M$ is called \emph{flat} if the functor
              \[
                  -\otimes_R M
              \]
              is exact.
    \end{enumerate}

    If $R$ is noncommutative, then in order for
    \[
        M\otimes_R -
    \]
    to be well-defined,$M$ must be a right $R$-module.
\end{definition}

